% ============================================
% DOC-16
% Gerado de molambudos.md — build_miolo.py
% ============================================
\section*{\textbf{DOC-16} --- Diário de Oliveira VI}
\markboth{DOC-16 --- Diário de Oliveira VI}{}

\noindent \textit{O último registro}

\noindent \textbf{Dr. Heitor Oliveira --- Caderno Pessoal}
\textbf{--- --- --- ---}

\vspace{0.5em}\noindent\rule{\textwidth}{0.5pt}\vspace{0.5em}

\noindent \textit{Este registro foi encontrado dentro do diário do paciente 1.260, em 2026, quando o diário foi devolvido ao Arquivo Público Mineiro. A caligrafia é de Oliveira, mas as últimas linhas alternam entre sua letra e a de Joaquim. A data está ilegível.}

\vspace{0.5em}\noindent\rule{\textwidth}{0.5pt}\vspace{0.5em}

\noindent Não sei que dia é.

\noindent Perdi a conta depois da vala.

\noindent Estou em Quixeramobim. Numa pensão. O quarto é pequeno, tem uma cama, um abajur, e o cheiro. O cheiro está em todo lugar agora. Na roupa. Na pele. Na comida. A comida tem gosto de doce. Tudo tem gosto de doce.

\noindent O olho direito está completamente amarelo.

\noindent Não olho no espelho.

\noindent O diário escreve sozinho agora. Eu não toco nele. Ele está sobre a mesa. As páginas viram sozinhas. A letra aparece. Eu leio.

\noindent Ele escreveu:

\noindent \textit{\textquotedblleft{}Heitor. Você está pronto.\textquotedblright{}}

\noindent Respondi: \textquotedblleft{}Pronto para quê?\textquotedblright{}

\noindent \textit{\textquotedblleft{}Para aceitar.\textquotedblright{}}

\noindent \textquotedblleft{}Aceitar o quê?\textquotedblright{}

\noindent \textit{\textquotedblleft{}Que você não é mais médico. Você é paciente. Paciente 1.261. O número não é de Joaquim. É da criatura. Ela conta os corpos que ocupou. Você é o próximo da fila.\textquotedblright{}}

\noindent Peguei uma caneta. Escrevi por cima:

\noindent \textit{\textquotedblleft{}Não aceito.\textquotedblright{}}

\noindent O diário respondeu:

\noindent \textit{\textquotedblleft{}Não precisa aceitar. Só precisa existir. A criatura vive em você. Aceitando ou não.\textquotedblright{}}

\vspace{0.5em}\noindent\rule{\textwidth}{0.5pt}\vspace{0.5em}

\noindent Hoje de manhã, uma mulher bateu na porta.

\noindent Era idosa. Cabelos brancos. Rosto cansado. Ela me olhou e disse:

\noindent \textit{\textquotedblleft{}O senhor é o médico?\textquotedblright{}}

\noindent \textit{\textquotedblleft{}...Era.\textquotedblright{}}

\noindent \textit{\textquotedblleft{}O senhor leu o diário?\textquotedblright{}}

\noindent \textit{\textquotedblleft{}Como sabe?\textquotedblright{}}

\noindent \textit{\textquotedblleft{}O olho. O amarelo. Todo mundo que leu fica com o olho.\textquotedblright{}}

\noindent \textit{\textquotedblleft{}Quem é a senhora?\textquotedblright{}}

\noindent Ela sentou na cama. Sem pedir.

\noindent \textit{\textquotedblleft{}Meu nome é Maria do Rosário. Sou mãe do Joaquim.\textquotedblright{}}

\noindent Minha boca abriu.

\noindent \textit{\textquotedblleft{}A senhora está viva?\textquotedblright{}}

\noindent \textit{\textquotedblleft{}Estava. Morri em 1965. Mas isso não importa. O que importa é que eu vim avisar o senhor.\textquotedblright{}}

\noindent \textit{\textquotedblleft{}Avisar o quê?\textquotedblright{}}

\noindent \textit{\textquotedblleft{}O diário não pode ficar com o senhor. O senhor tem que devolver.\textquotedblright{}}

\noindent \textit{\textquotedblleft{}Devolver para onde?\textquotedblright{}}

\noindent \textit{\textquotedblleft{}Para o arquivo. Para que outra pessoa leia.\textquotedblright{}}

\noindent \textit{\textquotedblleft{}Outra pessoa? Para contaminar mais alguém?\textquotedblright{}}

\noindent Ela me olhou. Os olhos dela estavam cansados.

\noindent \textit{\textquotedblleft{}O senhor acha que eu não tentei? Quando descobri que Joaquim estava vivo, em 1950, vim visitar. Ele me deu o diário. Disse: 'Leia, mãe. Para saber quem eu sou.' Eu li. Passei uma semana com ele. Depois devolvi. Disse que não queria saber.\textquotedblright{}}

\noindent \textit{\textquotedblleft{}A senhora sentiu o cheiro?\textquotedblright{}}

\noindent \textit{\textquotedblleft{}Senti.\textquotedblright{}}

\noindent \textit{\textquotedblleft{}O olho?\textquotedblright{}}

\noindent \textit{\textquotedblleft{}Clareou. Mas passou. Passei o diário adiante. Dei para o coveiro. Ele enterrou com Joaquim. Pensei que tinha acabado.\textquotedblright{}}

\noindent \textit{\textquotedblleft{}Não acabou.\textquotedblright{}}

\noindent \textit{\textquotedblleft{}Não. O coveiro leu. Depois devolveu ao arquivo. Depois veio o Dr. Tavares. Depois a Dra. Verônica. Depois o senhor.\textquotedblright{}}

\noindent Ela levantou.

\noindent \textit{\textquotedblleft{}O senhor não é o primeiro, doutor. Não vai ser o último.\textquotedblright{}}

\noindent \textit{\textquotedblleft{}O que eu faço?\textquotedblright{}}

\noindent Ela foi até a porta.

\noindent \textit{\textquotedblleft{}Passa adiante.\textquotedblright{}}

\noindent \textit{\textquotedblleft{}Para quem?\textquotedblright{}}

\noindent \textit{\textquotedblleft{}Para alguém que precise saber. Para alguém que possa contar a história.\textquotedblright{}}

\noindent \textit{\textquotedblleft{}E se eu não passar?\textquotedblright{}}

\noindent Ela olhou para mim. Pela última vez.

\noindent \textit{\textquotedblleft{}Então o senhor morre com ela dentro de você. E o ciclo espera. Até o próximo leitor.\textquotedblright{}}

\noindent Ela saiu.

\noindent Nunca mais a vi.

\vspace{0.5em}\noindent\rule{\textwidth}{0.5pt}\vspace{0.5em}

\noindent Hoje é --- não sei.

\noindent O diário está aberto. Não toquei nele. Mas ele escreveu:

\noindent \textit{\textquotedblleft{}Volte para Barbacena. Devolva o diário. Deixe que outro leia. É a única saída.\textquotedblright{}}

\noindent Vou fazer isso.

\noindent Vou devolver.

\noindent Mas antes, vou escrever tudo.

\noindent Para que quem ler depois de mim saiba.

\noindent \textbf{O ciclo não acaba.}
\textbf{O ciclo se repete.}
\textbf{O paciente 1.260 morreu.}
\textbf{O paciente 1.261 sou eu.}
\textbf{O paciente 1.262 será outro.}

\noindent \textit{--- Heitor Oliveira}
\textit{--- Paciente 1.261}
\textit{--- Ex-médico}
\textit{--- Ex-homem}
\textit{--- Ex-tudo}

\vspace{0.5em}
\noindent\textit{\textbf{↪ Rotas:} \rota{DOC-01} \rotasep \rota{CONT-05} \rotasep \rota{DOC-21}}